% GENERATED FILE. Edit canonical lessons, then run npm run generate:books.
% canonical-source-hash: fnv1a64:6f12b2e4222caf48
% canonical-lessons: KA-C525-bhava, KA-C525-nadini, KA-C525-maiduna, KA-C525-sodarasambandhi, KA-C525-sahodara, KA-R525-review-519-525

\chapter{In-laws and Cousins}
\label{ch:ka-in-laws-and-cousins}
\hlchaptermodality{525}

\begin{chapteropening}
\textbf{By the end of this chapter:} \emph{I can say a brother-in-law, husband's sister, husband's younger brother, a cousin and a brother.}

Complete the last lesson of chapter 525: I can say a brother-in-law, husband's sister, husband's younger brother, a cousin and a brother.
\end{chapteropening}

\section[\texorpdfstring{\kn{ಭಾವ} (bhāva)}{ಭಾವ (bhāva)}]{\kn{ಭಾವ} (bhāva) — a brother-in-law}

\label{lesson:KA-C525-bhava}



\begin{quote}
Before the new one: say the Kannada for mother's elder sister, then the Kannada for a nephew.
\end{quote}

\subsection*{You'll want to know: \kn{ಭಾವ}}
\textbf{\kn{ಭಾವ}} — \emph{bhāva} — \textquotedblleft{}a brother-in-law\textquotedblright{}.

\begin{quote}
\kn{ನನ್ನ} \kn{ಅಕ್ಕನ} \kn{ಗಂಡ}, \kn{ಅಂದರೆ} \kn{ನನ್ನ} \kn{ಭಾವ}, \kn{ವೈದ್ಯರು}.
\end{quote}

\begin{quote}
\emph{nanna akkana gaṇḍa, andare nanna bhāva, vaidyaru.}
\end{quote}

\begin{quote}
My elder sister's husband, that is, my brother-in-law, is a doctor.
\end{quote}

An elder sister's husband, or a husband's elder brother.

\begin{cousinweb}
\textbf{\kn{ಭಾವ}} (\emph{bhāva}), the kinship word, is usually traced to Sanskrit \emph{bhāvuka}, \textquotedblleft{}a sister's husband\textquotedblright{} (Monier-Williams). It is a different word from \emph{bhāva}, \textquotedblleft{}being, feeling\textquotedblright{}.
\end{cousinweb}

\begin{grammarlens}[title={What you've built}]
One more word for describing what you lived and felt.
\end{grammarlens}

\subsection*{Your turn}
\begin{itemize}
\raggedright
  \item \emph{Say it:} \emph{bhāva}
  \item \emph{Say it:} \emph{bhāva}, once more
  \item \emph{Say it:} \emph{nanna akkana gaṇḍa, andare nanna bhāva, vaidyaru.} — My elder sister's husband, that is, my brother-in-law, is a doctor.
  \item \emph{Recall:} say the Kannada for mother's elder sister, then the Kannada for a nephew, then say \emph{bhāva} again
\end{itemize}

\begin{tcolorbox}[breakable,colback=teal!4,colframe=teal!35!black,title={Before you move on}]
What does \kn{ಭಾವ} mean? (\textquotedblleft{}A brother-in-law\textquotedblright{}.) Say it once more.
\end{tcolorbox}
\section[\texorpdfstring{\kn{ನಾದಿನಿ} (nādini)}{ನಾದಿನಿ (nādini)}]{\kn{ನಾದಿನಿ} (nādini) — husband's sister}

\label{lesson:KA-C525-nadini}



\begin{quote}
Before the new one: say the Kannada for a nephew, then the Kannada for a brother-in-law.
\end{quote}

\subsection*{You'll want to know: \kn{ನಾದಿನಿ}}
\textbf{\kn{ನಾದಿನಿ}} — \emph{nādini} — \textquotedblleft{}husband's sister\textquotedblright{}.

\begin{quote}
\kn{ನನ್ನ} \kn{ನಾದಿನಿ} \kn{ನಮ್ಮ} \kn{ಮನೆಯ} \kn{ಹತ್ತಿರ} \kn{ಇರುತ್ತಾಳೆ}.
\end{quote}

\begin{quote}
\emph{nanna nādini namma maneya hattira iruttāḷe.}
\end{quote}

\begin{quote}
My husband's sister lives near our house.
\end{quote}

Also used for a wife's sister in some families.

\begin{cousinweb}
\textbf{\kn{ನಾದಿನಿ}} (\emph{nādini}) is native Dravidian: DEDR 3644 lists Kannada \emph{nādani}, \emph{nādini}, \textquotedblleft{}husband's sister, brother's wife\textquotedblright{}.
\end{cousinweb}

\begin{grammarlens}[title={What you've built}]
One more word for describing what you lived and felt.
\end{grammarlens}

\subsection*{Your turn}
\begin{itemize}
\raggedright
  \item \emph{Say it:} \emph{nādini}
  \item \emph{Say it:} \emph{nādini}, once more
  \item \emph{Say it:} \emph{nanna nādini namma maneya hattira iruttāḷe.} — My husband's sister lives near our house.
  \item \emph{Recall:} say the Kannada for a nephew, then the Kannada for a brother-in-law, then say \emph{nādini} again
\end{itemize}

\begin{tcolorbox}[breakable,colback=teal!4,colframe=teal!35!black,title={Before you move on}]
What does \kn{ನಾದಿನಿ} mean? (\textquotedblleft{}Husband's sister\textquotedblright{}.) Say it once more.
\end{tcolorbox}
\section[\texorpdfstring{\kn{ಮೈದುನ} (maiduna)}{ಮೈದುನ (maiduna)}]{\kn{ಮೈದುನ} (maiduna) — husband's younger brother}

\label{lesson:KA-C525-maiduna}



\begin{quote}
Before the new one: say the Kannada for a brother-in-law, then the Kannada for husband's sister.
\end{quote}

\subsection*{You'll want to know: \kn{ಮೈದುನ}}
\textbf{\kn{ಮೈದುನ}} — \emph{maiduna} — \textquotedblleft{}husband's younger brother\textquotedblright{}.

\begin{quote}
\kn{ನನ್ನ} \kn{ಮೈದುನ} \kn{ಎಂಜಿನಿಯರ್} \kn{ಆಗಿ} \kn{ಕೆಲಸ} \kn{ಮಾಡುತ್ತಾನೆ}.
\end{quote}

\begin{quote}
\emph{nanna maiduna eñjiniyar āgi kelasa māḍuttāne.}
\end{quote}

\begin{quote}
My husband's younger brother works as an engineer.
\end{quote}

Also a wife's brother in some families.

\begin{cousinweb}
\textbf{\kn{ಮೈದುನ}} (\emph{maiduna}) is a Kannada kinship word for a younger brother-in-law; its origin is not settled.
\end{cousinweb}

\begin{grammarlens}[title={What you've built}]
One more word for describing what you lived and felt.
\end{grammarlens}

\subsection*{Your turn}
\begin{itemize}
\raggedright
  \item \emph{Say it:} \emph{maiduna}
  \item \emph{Say it:} \emph{maiduna}, once more
  \item \emph{Say it:} \emph{nanna maiduna eñjiniyar āgi kelasa māḍuttāne.} — My husband's younger brother works as an engineer.
  \item \emph{Recall:} say the Kannada for a brother-in-law, then the Kannada for husband's sister, then say \emph{maiduna} again
\end{itemize}

\begin{tcolorbox}[breakable,colback=teal!4,colframe=teal!35!black,title={Before you move on}]
What does \kn{ಮೈದುನ} mean? (\textquotedblleft{}Husband's younger brother\textquotedblright{}.) Say it once more.
\end{tcolorbox}
\section[\texorpdfstring{\kn{ಸೋದರಸಂಬಂಧಿ} (sōdarasambandhi)}{ಸೋದರಸಂಬಂಧಿ (sōdarasambandhi)}]{\kn{ಸೋದರಸಂಬಂಧಿ} (sōdarasambandhi) — a cousin}

\label{lesson:KA-C525-sodarasambandhi}



\begin{quote}
Before the new one: say the Kannada for husband's sister, then the Kannada for husband's younger brother.
\end{quote}

\subsection*{You'll want to know: \kn{ಸೋದರಸಂಬಂಧಿ}}
\textbf{\kn{ಸೋದರಸಂಬಂಧಿ}} — \emph{sōdarasambandhi} — \textquotedblleft{}a cousin\textquotedblright{}.

\begin{quote}
\kn{ಮದುವೆಯಲ್ಲಿ} \kn{ನನ್ನ} \kn{ಎಲ್ಲಾ} \kn{ಸೋದರಸಂಬಂಧಿಗಳನ್ನು} \kn{ನೋಡಿದೆ}.
\end{quote}

\begin{quote}
\emph{maduveyalli nanna ellā sōdarasambandhigaḷannu nōḍide.}
\end{quote}

\begin{quote}
I saw all my cousins at the wedding.
\end{quote}

Plural: \kn{ಸೋದರಸಂಬಂಧಿಗಳು}, \emph{sōdarasambandhigaḷu}.

Everyday talk often names the exact tie instead: \kn{ಮಾವನ} \kn{ಮಗ}, an uncle's son.

\begin{cousinweb}
\textbf{\kn{ಸೋದರಸಂಬಂಧಿ}} (\emph{sōdarasambandhi}) is a learned compound built from Sanskrit \emph{sodara}, \textquotedblleft{}born of the same womb\textquotedblright{}, and \emph{sambandha}, \textquotedblleft{}close connection; a relative\textquotedblright{} (Monier-Williams).
\end{cousinweb}

\begin{grammarlens}[title={What you've built}]
One more word for describing what you lived and felt.
\end{grammarlens}

\subsection*{Your turn}
\begin{itemize}
\raggedright
  \item \emph{Say it:} \emph{sōdarasambandhi}
  \item \emph{Say it:} \emph{sōdarasambandhi}, once more
  \item \emph{Say it:} \emph{maduveyalli nanna ellā sōdarasambandhigaḷannu nōḍide.} — I saw all my cousins at the wedding.
  \item \emph{Recall:} say the Kannada for husband's sister, then the Kannada for husband's younger brother, then say \emph{sōdarasambandhi} again
\end{itemize}

\begin{tcolorbox}[breakable,colback=teal!4,colframe=teal!35!black,title={Before you move on}]
What does \kn{ಸೋದರಸಂಬಂಧಿ} mean? (\textquotedblleft{}A cousin\textquotedblright{}.) Say it once more.
\end{tcolorbox}
\section[\texorpdfstring{\kn{ಸಹೋದರ} (sahōdara)}{ಸಹೋದರ (sahōdara)}]{\kn{ಸಹೋದರ} (sahōdara) — a brother}

\label{lesson:KA-C525-sahodara}



\begin{quote}
Before the new one: say the Kannada for husband's younger brother, then the Kannada for a cousin.
\end{quote}

\subsection*{You'll want to know: \kn{ಸಹೋದರ}}
\textbf{\kn{ಸಹೋದರ}} — \emph{sahōdara} — \textquotedblleft{}a brother\textquotedblright{}.

\begin{quote}
\kn{ನನಗೆ} \kn{ಇಬ್ಬರು} \kn{ಸಹೋದರರು} \kn{ಇದ್ದಾರೆ}.
\end{quote}

\begin{quote}
\emph{nanage ibbaru sahōdararu iddāre.}
\end{quote}

\begin{quote}
I have two brothers.
\end{quote}

Plural: \kn{ಸಹೋದರರು}, \emph{sahōdararu}.

Formal for \kn{ಅಣ್ಣ} or \kn{ತಮ್ಮ}; a sister is \kn{ಸಹೋದರಿ}.

\begin{cousinweb}
\textbf{\kn{ಸಹೋದರ}} (\emph{sahōdara}) is borrowed from Sanskrit \emph{sahodara}, \textquotedblleft{}co-uterine, born of the same womb\textquotedblright{} (Monier-Williams): \emph{saha}, \textquotedblleft{}with\textquotedblright{}, and \emph{udara}, \textquotedblleft{}womb\textquotedblright{}.
\end{cousinweb}

\begin{grammarlens}[title={What you've built}]
One more word for describing what you lived and felt.
\end{grammarlens}

\subsection*{Your turn}
\begin{itemize}
\raggedright
  \item \emph{Say it:} \emph{sahōdara}
  \item \emph{Say it:} \emph{sahōdara}, once more
  \item \emph{Say it:} \emph{nanage ibbaru sahōdararu iddāre.} — I have two brothers.
  \item \emph{Recall:} say the Kannada for husband's younger brother, then the Kannada for a cousin, then say \emph{sahōdara} again
\end{itemize}

\begin{tcolorbox}[breakable,colback=teal!4,colframe=teal!35!black,title={Before you move on}]
What does \kn{ಸಹೋದರ} mean? (\textquotedblleft{}A brother\textquotedblright{}.) Say it once more.
\end{tcolorbox}
\section[(dialogue)]{Looking back: from Running a Business to In-laws and Cousins}

\label{lesson:KA-R525-review-519-525}



\begin{quote}
Say the Kannada for an entrepreneur and a manager.
\end{quote}

\subsection*{Your turn}
Say these aloud:

\begin{itemize}
\raggedright
  \item an entrepreneur, a manager, an assistant, to assign, to save up
  \item to get involved, to take charge of, to get practised, supervision, punctuality
  \item to toil, oversight, to insist, hard work, rivalry
  \item a blacksmith, a goldsmith, a weaver, a sculptor, to carve
  \item family life, a homemaker, a partner, the pain of parting, married life
  \item father's younger brother, father's elder brother, mother's younger sister, mother's elder sister, a nephew
  \item a brother-in-law, husband's sister, husband's younger brother, a cousin, a brother
\end{itemize}

\begin{tcolorbox}[breakable,colback=teal!4,colframe=teal!35!black,title={Before you move on}]
An entrepreneur? (\textbf{\kn{ಉದ್ಯಮಿ}}, \emph{udyami}.) A manager? (\textbf{\kn{ವ್ಯವಸ್ಥಾಪಕ}}, \emph{vyavasthāpaka}.) An assistant? (\textbf{\kn{ಸಹಾಯಕ}}, \emph{sahāyaka}.) To assign? (\textbf{\kn{ನಿಯೋಜಿಸು}}, \emph{niyōjisu}.) To save up? (\textbf{\kn{ಕೂಡಿಡು}}, \emph{kūḍiḍu}.) To get involved? (\textbf{\kn{ತೊಡಗು}}, \emph{toḍagu}.) To take charge of? (\textbf{\kn{ವಹಿಸು}}, \emph{vahisu}.) To get practised? (\textbf{\kn{ಪಳಗು}}, \emph{paḷagu}.) Supervision? (\textbf{\kn{ಮೇಲ್ವಿಚಾರಣೆ}}, \emph{mēlvicāraṇe}.) Punctuality? (\textbf{\kn{ಸಮಯಪಾಲನೆ}}, \emph{samayapālane}.) To toil? (\textbf{\kn{ದುಡಿ}}, \emph{duḍi}.) Oversight? (\textbf{\kn{ಉಸ್ತುವಾರಿ}}, \emph{ustuvāri}.) To insist? (\textbf{\kn{ಒತ್ತಾಯಿಸು}}, \emph{ottāyisu}.) Hard work? (\textbf{\kn{ಪರಿಶ್ರಮ}}, \emph{pariśrama}.) Rivalry? (\textbf{\kn{ಪೈಪೋಟಿ}}, \emph{paipōṭi}.) A blacksmith? (\textbf{\kn{ಕಮ್ಮಾರ}}, \emph{kammāra}.) A goldsmith? (\textbf{\kn{ಅಕ್ಕಸಾಲಿಗ}}, \emph{akkasāliga}.) A weaver? (\textbf{\kn{ನೇಕಾರ}}, \emph{nēkāra}.) A sculptor? (\textbf{\kn{ಶಿಲ್ಪಿ}}, \emph{śilpi}.) To carve? (\textbf{\kn{ಕೆತ್ತು}}, \emph{kettu}.) Family life? (\textbf{\kn{ಸಂಸಾರ}}, \emph{saṃsāra}.) A homemaker? (\textbf{\kn{ಗೃಹಿಣಿ}}, \emph{gṛhiṇi}.) A partner? (\textbf{\kn{ಸಂಗಾತಿ}}, \emph{saṅgāti}.) The pain of parting? (\textbf{\kn{ವಿರಹ}}, \emph{viraha}.) Married life? (\textbf{\kn{ದಾಂಪತ್ಯ}}, \emph{dāmpatya}.) Father's younger brother? (\textbf{\kn{ಚಿಕ್ಕಪ್ಪ}}, \emph{cikkappa}.) Father's elder brother? (\textbf{\kn{ದೊಡ್ಡಪ್ಪ}}, \emph{doḍḍappa}.) Mother's younger sister? (\textbf{\kn{ಚಿಕ್ಕಮ್ಮ}}, \emph{cikkamma}.) Mother's elder sister? (\textbf{\kn{ದೊಡ್ಡಮ್ಮ}}, \emph{doḍḍamma}.) A nephew? (\textbf{\kn{ಸೋದರಳಿಯ}}, \emph{sōdaraḷiya}.) A brother-in-law? (\textbf{\kn{ಭಾವ}}, \emph{bhāva}.) Husband's sister? (\textbf{\kn{ನಾದಿನಿ}}, \emph{nādini}.) Husband's younger brother? (\textbf{\kn{ಮೈದುನ}}, \emph{maiduna}.) A cousin? (\textbf{\kn{ಸೋದರಸಂಬಂಧಿ}}, \emph{sōdarasambandhi}.) A brother? (\textbf{\kn{ಸಹೋದರ}}, \emph{sahōdara}.)
\end{tcolorbox}
